\chapter{绪论}
\label{chap:intro}

\section{研究背景与意义}
\label{sec:bg}

近年来，随着物联网（Internet of Things, IoT）技术迅速发展，物联网设备已被广泛用于智能家居、智慧城市、新能源、车联网、医疗健康等多个应用领域。根据 GSMA 发布的《The mobile economy 2020》报告显示，2019 年全球物联网总连接数达到 120 亿，预计到 2025 年，全球物联网总连接数规模将达到 246 亿，年复合增长率高达 13\%\cite{gsma2020mobile}。物联网操作系统是一种面向物联网设备和应用的轻量级嵌入式操作系统，它可以提供更好的设备管理、数据处理和安全性能。目前，市场上有许多物联网操作系统，例如华为的 OpenHarmony\cite{openharmony}、亚马逊的 FreeRTOS\cite{freertos}、Linux 基金会的 Zephyr\cite{zephyr} 等。然而，物联网和嵌入式操作系统面临着严峻的安全挑战\cite{yang2021iotsurvey}，如恶意攻击、数据泄露、隐私侵犯等，这些安全问题可能会给用户带来损失和危害。在真实世界中发生的物联网安全问题也层出不穷。2018 年，一款名为"Mirai"的恶意软件\cite{mirai2018}利用了 IoT 设备的弱点，通过攻击嵌入式设备的操作系统，将其变成僵尸网络的一部分，对互联网上的服务器发起了大规模的 DDoS 攻击。俄乌冲突期间，乌克兰遭遇了俄罗斯发起的破坏性网络攻击，该攻击利用一种名为 Industroyer 的恶意软件控制电力系统中的开关和断路器，意图瘫痪乌克兰的电力系统。这些安全事件的不断发生表明，当前物联网和嵌入式操作系统安全问题的研究刻不容缓。

固件仿真（Firmware Re-Hosting）技术是一种克服固件分析硬件依赖问题的有效方法，它通过在商用计算机和服务器的软件模拟器上运行嵌入式系统固件，实现对嵌入式设备的软件级模拟，使得包含嵌入式系统的固件无需实际硬件即可进行安全测试\cite{firmadyne}。然而，当前固件仿真技术的\textbf{自动化程度仍然严重不足}，这主要体现在以下几个方面：首先，现有的外设仿真方法大多需要人工编写模型或基于有限规则的自动推断。硬件在环方案需要为每种外设手动配置物理设备连接；基于原始硬件的外设建模方案需要人工分析数据手册并编写仿真代码；基于程序分析的方案虽然能自动推断部分外设行为，但推断规则有限，难以覆盖复杂外设的状态机逻辑。其次，驱动层替换方案（如 HALucinator）虽然利用了微控制器制造商提供的硬件抽象层（HAL），通过函数 Hook 的方式进行替换，但驱动函数的定位、语义理解和替换代码生成仍然需要大量人工干预。分析人员需要手动识别固件中哪些部分存在硬件依赖并需要替换，理解驱动函数与硬件的交互逻辑，并编写相应的替代代码，这个过程不仅耗时耗力，而且容易出错。这种\textbf{自动化程度不足导致的人工成本高昂}，严重制约了固件仿真技术在大规模安全测试中的应用，成为当前物联网设备安全领域亟待解决的关键问题。

现有固件仿真技术的自动化不足主要体现在三个核心环节。在驱动函数定位环节，识别固件中哪些部分存在硬件依赖并需要替换是一项具有挑战性的任务。在嵌入式程序中，硬件依赖代码往往分散在各个模块中，通过函数调用、全局变量访问、寄存器操作等多种形式存在，手动定位这些代码需要分析人员对固件结构和硬件特性有深入理解，过程既耗时又容易出错。在外设建模环节，当前的自动化方法只能处理简单的外设行为，对于网络控制器、USB 设备等具有多通道、DMA 传输、中断嵌套等特性的复杂外设，往往需要人工编写仿真模型。在替换代码生成环节，现有工具主要通过函数 Hook 返回预设值或执行简化逻辑，缺乏对驱动函数内部语义的深入理解，无法准确模拟硬件行为，需要人工编写复杂的替代代码。这三个环节的自动化不足相互叠加，使得固件仿真过程对分析人员经验的依赖程度过高，难以适应大规模、多样化的固件安全测试需求。

针对上述问题，本研究致力于提升固件仿真技术的自动化程度，降低人工成本，实现高效的物联网设备固件安全测试。具体而言，本研究提出一种基于驱动程序替换的固件仿真方法，通过静态分析技术自动识别固件中的硬件依赖代码，利用大语言模型的代码理解能力深入分析驱动函数语义并自动生成可替代的驱动实现，构建动态测试反馈闭环以迭代优化仿真质量。该方法的核心创新在于将自动化技术贯穿于驱动函数定位、语义分析和代码生成的全过程，显著减少了对分析人员人工干预的依赖。通过提升固件仿真的自动化水平，本研究不仅能够降低固件安全测试的成本，还能够提高测试效率和覆盖率，为物联网设备的大规模安全测试提供技术支撑，具有重要的理论意义和应用价值。

\section{国内外研究现状}
\label{sec:related}

固件是嵌入式设备中的软件组成部分，通常包含了设备的操作系统、驱动程序、应用程序和配置文件等。固件的安全性对于保护设备和用户数据的完整性和隐私至关重要。然而，由于固件的复杂性、多样性和闭源性，对固件进行有效的安全测试是一项具有挑战性的任务。

基于固件仿真的模糊测试技术在近年来受到了越来越多的关注和研究，出现了一些代表性的工具和系统，如 FIRMADYNE\cite{firmadyne}、FUZZWARE\cite{fuzzware} 等。这些工具和系统在不同程度上实现了对固件仿真和模糊测试的支持，并在实际的固件中发现了一些新的漏洞。然而，基于固件仿真的模糊测试技术仍然存在一些挑战和问题，如仿真环境的准确性、测试用例的有效性等，需要进一步的改进和优化。下面将从固件仿真技术、嵌入式设备模糊测试技术两个方面阐述当前基于固件仿真的嵌入式设备固件模糊测试技术研究现状。

\subsection{固件仿真技术研究现状}
\label{sec:firmware_sim}

当前固件仿真的主要挑战是为模拟器中没有实现的硬件提供有效的输入。诸如 QEMU\cite{qemu} 之类的仿真器提供了模拟各种 CPU 体系结构的能力，但对商用嵌入式系统中使用的各种各样的外设，例如计时器、UART、以太网控制器等，提供的支持却十分有限，而这些外设对于正确实现固件仿真却是十分必要的。

固件仿真技术的发展经历了从依赖真实硬件到全自动化模拟的演进过程。早期研究采用硬件在环（Hardware in the Loop, HIL）的方法，如 Zaddich 等人提出的 Avatar\cite{avatar} 和 Muench 等人提出的 Avatar2\cite{avatar2}，通过混合执行的方式将固件分析过程中对外设的访问转发到物理硬件中。然而，这种方案仍然依赖实际设备，且仿真环境和物理设备之间的切换影响了运行速度。为摆脱对真实硬件的依赖，研究者提出了多种全自动化仿真方案，主要包括以下三类。

\subsubsection{基于原始硬件的外设建模方案}

解决外设依赖的一种方法是在底层内存映射寄存器接口上实现硬件，或通过观察真实硬件的行为来构建外设模型。然而这是一项费力的任务，扩展性不好。当前主流的仿真器 QEMU 也仅能仿真 10 余种设备，远不能满足固件仿真工作的需求。

Gustafson 等人提出的 Pretender\cite{pretender} 工具通过观察固件与物理设备之间真实发生的交互记录，使用机器学习、模式识别等方法为每个外设构建模型，再结合 QEMU 或 angr 等仿真引擎实现固件仿真。这种方法需要真实设备的支持，但在观测到足够的交互模式后，可以自动生成较为准确的外设模型。

\subsubsection{基于程序分析的外设建模方案}

基于程序分析的外设建模方案通过对固件代码进行静态或动态分析，提取固件与外设之间的交互逻辑，构建虚拟外设模型以替代真实硬件。

Feng 等人提出的 P2IM\cite{fengp2020p2im} 工具通过动态符号执行和模糊测试，自动识别固件中的输入点和输出点，构建与真实设备相似的虚拟模型。Cao 等人提出的 Laelaps\cite{laelaps} 在因访问未实现的外设而卡住时，切换到符号执行以找到正确输入，引导 QEMU 找到与实际执行最相近的路径。Zhou 等人提出的 uEmu\cite{uemu} 工具利用符号执行分析固件镜像，将未知的外设寄存器表示为符号，构建推断规则并存储在知识库中。Scharnowski 等人提出的 Fuzzware\cite{fuzzware} 使用精确的 MMIO 建模实现有效的固件模糊测试，通过确定固件逻辑如何实际使用硬件生成的值，自动生成模型以提高模糊测试输入数据的有效性。

\subsubsection{基于系统或驱动层替换的外设仿真方案}

一种有前途的方法是基于系统或驱动层替换的仿真技术，或称为高级仿真（High Level Emulation, HLE）技术。它利用固件中的公共抽象来消除为外设提供低级支持的需要。

Chen 等人提出了一个针对 Linux 内核的物联网设备固件仿真平台 Firmadyne\cite{firmadyne}，通过修改内核和驱动来添加硬件支持，实现物联网设备的全系统仿真。该工具解决了定制化外围 I/O、非易失性内存、动态文件等仿真技术难题。Li 等人创建了 Para-rehosting\cite{pararehosting} 工具，该工具的主要思想是利用 POSIX 接口抽象并实现了一个便携式 MCU，以将嵌入式系统源代码重新编译到商用机器的本机可执行文件中，以便利用 x86 上许多现有的复杂软件测试工具对代码进行测试。然而，这种方案的实现需要分析师获得固件样本的源代码，而在第三方的固件分析工作中，源代码的获取往往是困难甚至完全不可行的。Clements 等人提出了 HALucinator\cite{clements2020halucinator} 工具，该工具利用了微控制器制造商提供的硬件抽象层（Hardware Abstraction Layer, HAL），通过对硬件抽象层函数的抽象与替换进行固件仿真，这项工作在重托管简单的裸金属固件程序时取得了很好的效果。

然而，上述基于系统或驱动层替换的仿真技术自动化程度较差，难以适应大规模安全测试的需要。具体而言，这类技术面临以下两个核心难点：一是难以自动化定位需要进行替换的驱动程序。在嵌入式程序中，识别哪些部分存在硬件依赖并需要替换，往往依赖分析人员的手动干预。手动分析嵌入式程序的复杂性使得这一过程既耗时又容易出错。二是难以自动化生成替换程序。除了识别硬件相关功能，当前的技术也缺乏自动化生成替代程序的能力。分析人员必须手动分析函数的功能，理解其与硬件和中间件的交互，并编写相应的替代代码。这个过程不仅繁琐，而且容易导致人力成本的增加。

\subsubsection{固件仿真技术对比}

表~\ref{tab:sim_comparison} 对比了当前主流固件仿真技术的优缺点。

\begin{table}[htbp]
\centering
\caption{当前主流固件仿真技术对比}
\label{tab:sim_comparison}
\begin{tabular}{p{2.2cm}p{2.5cm}p{4cm}p{4cm}}
\toprule
\textbf{方案} & \textbf{典型工具} & \textbf{优点} & \textbf{缺点} \\
\midrule
基于原始硬件的外设建模 & QEMU, Pretender & 可对外设行为进行精确模拟，提高测试效率和覆盖率 & 需要为每种外设专门建模，工作量大，通用性差，依赖真实设备观测 \\
\midrule
基于程序分析的外设建模 & P2IM, uEmu, Laelaps, Fuzzware & 可根据程序自身信息自动推断外设行为，减少人工干预，提高通用性 & 难以完全覆盖复杂外设的所有行为，可能存在误报或漏报 \\
\midrule
基于系统或驱动层替换 & HALucinator, Firmadyne, Para-rehosting & 在系统层或驱动层对函数进行抽象和替换，降低对硬件的依赖，提高测试速度和灵活性 & 驱动函数的定位与替换依赖人工干预，自动化程度不足，难以适应大规模安全测试 \\
\bottomrule
\end{tabular}
\end{table}

在这三类方案中，基于系统或驱动层替换的固件仿真方法是一种有前景的方案。它从固件系统层和驱动函数入手，较好地解决了基于程序分析和原始硬件建模的方法对复杂外设难以完全模拟或建模工作量大等痛点。HALucinator 等工具利用微控制器制造商提供的硬件抽象层（Hardware Abstraction Layer, HAL），通过函数替换的方式实现固件仿真，在重托管简单的裸金属固件程序时取得了很好的效果。

然而，现有驱动层替换方案在自动化方面仍存在明显不足：一方面，驱动函数的定位依赖分析人员的手动干预，识别固件中哪些部分存在硬件依赖并需要替换是一个耗时且容易出错的过程；另一方面，替换程序的生成也缺乏自动化支持，分析人员需要手动分析函数功能、理解其与硬件的交互，并编写相应的替代代码。为解决上述问题，需要设计一种能够自动化定位需要替换的驱动函数、并利用大模型或代码分析技术自动生成替代程序的方案，以提升固件仿真过程的自动化程度，实现高效的物联网设备程序仿真与动态测试。

\subsection{嵌入式系统模糊测试技术发展现状}
\label{sec:embedded_fuzz}

嵌入式系统模糊测试技术是一种针对嵌入式设备的软件安全测试方法，它可以自动或半自动地生成随机或变异的输入数据，发送给嵌入式设备的 Web 服务或其他接口，监控设备的运行状态，发现潜在的漏洞或缺陷。目前，嵌入式系统模糊测试技术主要从以下三个方向进行发展：

\subsubsection{面向真实设备的黑盒模糊测试}

直接对嵌入式设备进行模糊测试，不需要对设备进行任何修改或插桩，只需要获取设备的网络接口或串口等信息，然后使用模糊测试工具生成并发送测试用例，监控设备的响应和异常。Chen 等人开发的 IoTFuzzer\cite{iotfuzzer} 工具，就利用了应用程序与物联网设备的交互信息，通过网络接口向嵌入式设备发送模糊测试数据，进行模糊测试工作。这种方法的优点是简单易用，不依赖于设备的内部结构或代码，适用于黑盒测试场景。缺点是测试效率较低，覆盖率较差，难以发现深层次的漏洞。

\subsubsection{基于固件仿真的灰盒模糊测试}

基于固件仿真的模糊测试，通过对嵌入式设备的固件进行逆向工程，提取出可执行的二进制代码，然后使用仿真器在虚拟环境中运行固件，并对其进行模糊测试。基于该方案的仿真工具有 Fuzzware\cite{fuzzware}、uEmu\cite{uemu}、HALucinator\cite{clements2020halucinator} 等。这种方法的优点是可以对固件进行插桩或修改，获取更多的运行信息，提高测试效率和覆盖率，发现更多的漏洞。缺点是需要对固件进行复杂的逆向分析，可能遇到加密或保护机制，同时仿真器可能存在不完全或不准确的问题，导致测试结果与真实设备不一致。

\subsection{大模型程序分析与代码生成技术研究现状}
\label{sec:llm_related}

大语言模型（Large Language Model, LLM）在代码理解和生成方面展现出了强大的能力。以 GPT-4\cite{openai2023gpt4}、Code Llama\cite{roziere2023code}、StarCoder\cite{li2023starcoder} 等为代表的代码大模型能够理解复杂的代码逻辑，并生成高质量的代码。近年来，研究者开始探索将 LLM 应用于程序分析和代码生成领域，涌现出多种增强策略以提升模型在复杂编程任务中的表现。

\subsubsection{面向特定任务的代码生成方法}

早期 LLM 通过零样本提示能够生成基础代码，但在处理复杂编程任务时准确性和可靠性仍存在不足。为克服这些局限，研究者发展了以下增强策略：

\textbf{指令微调（Instruction Tuning）}：通过对预训练 LLM 在特定领域数据集上进行精细调整，显著提升模型对特定编程任务的理解和执行能力\cite{wei2022finetuned}。DolphCoder 模型通过多样化的指令和自评估机制，结合代码评估目标有效提升了代码生成能力\cite{dolphcoder2024}。CodeLSI 框架则通过低秩优化与领域特定指令微调相结合，提高了代码生成的领域特异性和成本效益\cite{codelsi2024}。

\textbf{测试驱动的迭代生成}：利用测试用例作为反馈信号指导模型逐步修正生成的代码，直至通过所有测试。这种方法通过隐式执行监督信号，结合搜索、重排序或迭代修正等机制动态优化模型输出\cite{chen2023teaching}。RLCF（Reinforcement Learning with Compiler Feedback）方法通过编译器反馈对 LLM 进行额外训练，解决生成代码中可能存在的语言级不变量违规问题\cite{rlcf2024}。CodeDPO 框架则将偏好学习整合到代码生成过程中，优化模型优先选择正确且高效的解决方案\cite{codedpo2024}。

\textbf{多步推理方法}：思维链（Chain-of-Thought, CoT）和程序思维（Program-of-Thought, PoT）等方法通过将复杂问题分解为多个可管理的子步骤，提升模型解决复杂问题的能力\cite{wei2022cot}。MoTCoder 引入模块化思维指令框架来应对更具挑战性的编程任务\cite{motcoder2024}。程序思维方法则将中间推理步骤编译为可执行代码片段，实现了"推理即执行"，大幅提升了数值计算和 API 调用类任务的准确性。

\textbf{工具增强的大模型}：通过集成编译器、调试器、静态分析工具等外部工具，扩展模型的实际应用能力。这种范式将 LLM 定位为"AI 协调者"，通过标准化工具调用协议调度各种工具，形成感知-决策-执行的闭环\cite{schick2023toolformer}。MCCoder 系统专注于生成运动控制代码，并结合严格的验证机制以确保代码安全性\cite{mccoder2024}。RGD 代理调试器利用多 LLM 通过细化和生成指导来自动修复代码中的错误\cite{rgd2024}。

\subsubsection{系统级软件代码生成的挑战与前沿}

在系统级软件应用层面，特别是针对物联网设备驱动程序、操作系统内核和嵌入式固件等复杂系统，LLM 的代码生成面临独特挑战。IoT 设备驱动程序开发具有强硬件耦合性、并发中断敏感性、底层资源约束性以及隐式的 DMA/寄存器交互等特征，对 LLM 的代码生成能力提出了更高要求\cite{gpiot2024}。

当前前沿研究强调构建"分析—生成—验证"的闭环系统以应对这些挑战：

\textbf{静态分析集成}：静态分析工具如 CIL、Smatch 和 Frama-C 用于提取接口契约与控制流约束，对代码进行结构化表示\cite{spec2code2024}。这些工具能够识别不安全的功能、缓冲区溢出和竞态条件等安全漏洞，为 LLM 生成提供结构化上下文。

\textbf{动态符号执行验证}：动态符号执行引擎如 QEMU 和 KLEE 能够生成高覆盖率的路径约束，并反向验证生成代码的内存安全与中断语义正确性\cite{improvingllm2024}。通过结合编译器诊断、安全扫描和符号执行，LLM 可以迭代细化其输出。

\textbf{端到端流水线}：研究者正在构建端到端的"分析→生成→符号验证→反馈精炼"流水线\cite{agents4plc2024}。形式化规约被用作生成约束，而验证失败轨迹则被结构化为自然语言反馈以驱动下一轮生成。Agents4PLC 框架通过基于 LLM 的代理自动化了可编程逻辑控制器（PLC）代码生成与验证，提高了工业控制系统的操作效率和安全性。

这些研究为本文利用 LLM 进行驱动函数分析和替换提供了理论基础，特别是在构建自动化分析-生成-验证闭环方面具有重要的参考价值。

\subsection{现有研究的局限性分析}
\label{sec:limitations}

综合上述研究现状，现有固件仿真技术主要存在以下局限性：

\textbf{（1）外设建模依赖人工，难以处理复杂外设行为}

现有仿真方法大多依赖人工编写外设模型或基于有限规则的自动推断。硬件在环方案需要为每种外设手动配置物理设备连接；基于原始硬件的外设建模方案需要人工分析数据手册并编写仿真代码；基于程序分析的方案虽然能自动推断部分外设行为，但推断规则有限，难以覆盖复杂外设的状态机逻辑。例如，网络控制器、USB 设备等复杂外设具有多通道、DMA 传输、中断嵌套等特性，现有自动推断方法往往只能处理简单的 MMIO 寄存器读写，无法准确模拟这些复杂交互行为。

\textbf{（2）缺乏对驱动函数语义的深入理解，仿真结果可能不准确}

现有驱动层替换方案（如 HALucinator）主要通过函数 Hook 的方式进行替换，即当固件调用某个 HAL 函数时，返回预设的模拟值或执行简化的模拟逻辑。然而，这种方法缺乏对驱动函数内部语义的深入理解：函数之间的调用关系、参数的数据流传递、外设状态的变化等关键信息往往被忽略。这导致在某些场景下，仿真结果与真实硬件行为存在偏差，可能遗漏重要的程序路径或引入虚假的漏洞。

\textbf{（3）缺乏自动化反馈机制，难以迭代优化仿真质量}

当前大多数仿真工具采用"一次性配置"的模式：分析人员手动配置仿真环境后运行测试，如果发现问题则需要人工介入调试和修复。这种模式缺乏自动化的反馈闭环机制——当仿真过程中出现异常（如程序卡死、断言失败、外设访问超时）时，系统无法自动分析原因并调整仿真策略。这导致仿真质量的提升高度依赖分析人员的经验，难以适应大规模、多样化的固件测试需求。

\section{研究内容与主要贡献}
\label{sec:contribution}

为解决上述问题，本文提出一种基于驱动程序替换的物联网设备固件仿真方法。该方法通过静态分析自动识别固件中的硬件依赖代码，利用大语言模型分析驱动函数语义并自动生成可替代的驱动实现，构建动态测试反馈闭环以迭代优化仿真质量。本文主要的创新点和贡献有三个：

（1）本研究设计并实现了一个面向嵌入式固件的硬件依赖代码自动识别模块。该模块通过 MMIO 模式识别、数据流追踪与递归归纳等技术，从固件二进制代码中自动提取硬件相关的函数调用点、外设寄存器访问模式和驱动函数边界，构建可查询的驱动知识库。与现有依赖人工标注或有限规则推断的方法相比，该模块能够自动处理多种嵌入式平台（STM32、NXP 等）和操作系统（FreeRTOS、RT-Thread、裸机）的固件，显著降低了驱动函数定位的人工成本。

（2）本研究设计并实现了一个基于大语言模型的驱动函数分析与替换模块。该模块利用大语言模型的代码理解能力，对识别出的驱动函数进行语义分析，理解其与硬件的交互逻辑，并自动生成可在仿真环境中运行的替代实现。通过设计结构化的提示词模板和函数分类策略，本模块能够处理不同类型的驱动函数（如初始化、数据传输、中断处理等），并通过编译修复循环保障生成代码的可编译性。与现有基于预设返回值的简单 Hook 方法相比，该方法生成的替代代码具有更丰富的语义信息，能够更准确地模拟硬件行为。

（3）本研究引入了动态测试反馈闭环机制。当固件在仿真环境中运行时，该机制能够捕获运行时异常（如程序卡死、外设访问超时、断言失败等），分析异常原因，并将分析结果反馈给驱动替换模块以迭代优化替代代码。这种"生成-测试-反馈-修正"的闭环设计使仿真质量能够持续提升，减少了对分析人员经验的依赖。实验结果表明，该反馈机制能够在有限的迭代轮次内显著提高固件在仿真环境中的稳定性和运行覆盖率。

\section{论文组织结构}
\label{sec:structure}

本文共分为五章，组织结构如下：

\textbf{第一章：绪论}。介绍研究背景、意义、国内外研究现状以及本文的主要贡献。

\textbf{第二章：背景知识与相关技术}。介绍物联网设备硬件基础、嵌入式操作系统与驱动框架、大模型代码分析技术以及固件仿真与模糊测试基础。

\textbf{第三章：系统设计与实现}。详细描述系统的整体架构设计、各核心模块的设计与实现。

\textbf{第四章：实验设计与结果分析}。介绍实验环境、测试案例、实验结果及分析。

\textbf{第五章：总结与展望}。总结本文工作，讨论局限性，并展望未来研究方向。
